% critique.tex — f170v, vv. 679–720 : commentaire critique
% Discussions textuelles extraites du commentaire de l'auteur (commentaire.tex), texte et notes conservés ;
% répartition et retouches : repartition.md
\cc{702}{ὡς ἅπαν γὰρ ἐστὶν τόλμης ἔργα κἀναισχυντίας.}{Le début du vers dans le \italique{Ravennas}, ὡς ἅπαν γὰρ ... ἔργα κτλ., est fautif. La correction la plus simple à apporter est de lire : ὡς ἅπαντ᾽ ἄρ᾽ κτλ. Il n'y aurait alors qu'une simple confusion d'un Τ lu en Γ : ΑΠΑΝΤΑΡ $ \to $ ΑΠΑΝΓΑΡ.
Cette correction mineure fait concorder le texte d'Aristophane avec les parémiographes. Les variantes ne diffèrent plus que sur un seul mot : ἔργα (dans le \italique{Ravennas}) ; πλέα (\italique{Souda}) ; δοῦλα (Laurent.) ; $\emptyset$ (\italique{Corpus paroemiographorum}).

Rejetons d'emblée la leçon du \italique{Ravennas}, qui sent son interpolation ; ἔργα n'est qu'une glose d'ἅπαντα qui aura pénétré le texte.
La même erreur se produit au vers 717 des \italique{Oiseaux} : les manuscrits comportent `πρὸς ἅπαντα τρέπεσθε' ; une glose du scholiaste substitue ἔργα à ἅπαντα et Suidas (α 1638) finit par citer le vers ainsi : `πρὸς ἔργα τράπεσθε'.

La leçon de la \italique{Souda}, πλέα, est satisfaisante quant au sens (tout \italique{plein} d'effronterie et d'audace) ; elle n'est cependant pas métriquement correcte, raison pour laquelle Porson lui substitua μεστά. 
À défaut d'explication satisfaisante, j'ai préféré garder la leçon fautive de Suidas, non sans la frapper de l'obèle du désespoir (†).}
\cc{704}{οἷον ὑμῶν ἐξάρξω τὴν ἄγαν αὐθαδίαν.}{Il manque une syllabe à ἐξάρξω, la leçon du manuscrit de Ravenne. Les différents éditeurs y suppléent avec le verbe ἐξαράττειν (ἐξαράξω au futur ou, en s'éloignant davantage de la leçon -- quitte à corriger le texte... --, ἐξαράξει, pour faire de οἷον le sujet). La correction est légère et pourtant elle ne manque pas de laisser dubitatif : si l'image de \italique{briser} (ἐξαράξω) l'\italique{audace} (αὐθαδίαν) de \italique{quelqu'un} (ὑμῶν) ne paraît aucunement étrangère à une imagination française, il n'en allait peut-être pas de même pour les Grecs. L'incontournable Taillardat\nc{§628 de ses \italique{Images d'Aristophane}, ouvrage dont on ne soulignera jamais assez l'importance pour toute personne désireuse d'avoir un avant-goût du génie de la langue aristophanesque.} l'avait déjà noté : considérer l'arrogance «~comme un objet matériel qu'on peut briser [...] est une image unique~». En revanche, en grec comme en français, l'image de \italique{chasser}, \italique{pousser} hors (ἐξ-), \italique{repousser} un sentiment ou un trait de caractère est courante. Ainsi, l'on chassera l'audace par l'audace\nc{Athénée, \italique{Deipnosophistes}, II, 44a.} :\citComm{Ἀντιφάνης δέ· οἴνῳ .... τὸν οἶνον ἐξελαύνειν, [...] τριωβόλῳ δὲ πόρνην, αὐθαδίαν αὐθαδίᾳ.}L'on se débarrassera de son chagrin\nc{Aristote, \italique{Éthique à Nicomaque}, 1154a27.} :\citComm{ἐκκρούει τὴν λύπην}comme les vieillards cagneux secouent leurs ennuis chez Aristophane\nc{\italique{Grenouilles}, v. 346.} :\citComm{ἀποσείονται δὲ λύπας.}Euripide met dans la bouche du fourbe Polydectès qui veut s'accaparer Danaé que\nc{Fragment du \italique{Dictys}. Stobée, IV, 26, 16 $=$ fr. 7 \italique{Dictys} Jouan \& van Looy $=$ fr. 339 Kannicht $=$ fr. 433 Mette.} :\citComm{πατέρα τε παισὶν ἡδέως συνεκφέρειν\\
φίλους ἔρωτας \textbf{ἐκβαλόντ' αὐθαδίαν}.\\
Un père doit supporter avec indulgence les affaires amoureuses de son fils et \italique{repousser l'entêtement}.}
Le verbe ἀρήγω (futur : ἀρήξω) signifie \italique{secourir} et, par suite, \italique{écarter}, \italique{repousser}. Il conviendrait tant pour le sens que paléographiquement s'il était préfixé de ἐξ-.
Comme pour nous conforter, Hésychius (ε 3599) nous atteste justement l'existence du composé :\citComm{ἐξαρῆξαι· ἐκφορῆσαι.}Lisons donc :\citComm{ὑμῶν ἐξαρήξω τὴν ἄγαν αὐθαδίαν\\
Je vous extirperai votre trop-plein d'audace.}
P.-S. : l'idée de recourir aux matrices métaphoriques pour le bon établissement d'un texte revient à l'immense Jean Taillardat. Je ne peux que renvoyer à son coruscant article « Images et matrices métaphoriques »\nc{\italique{Bulletin de l'Association Guillaume Budé} (décembre 1977), Lettre d'humanité n° 36, pp. 344-354.}. Il le termine par ce vœu :\citComm{En guise de conclusion, je formulerai un souhait, celui qu'on puisse, un jour, consulter un dictionnaire des matrices métaphoriques pour les langues indo-européennes vivantes (langues populaires et argots compris) et mortes. On l'imaginerait volontiers comme une sorte de dictionnaire analogique, mais réduit aux images, les notions à exprimer (par exemple « colère », « inconsistance », « jeunesse », etc.) étant classées en ordre alphabétique ; quel merveilleux outil serait ainsi offert à qui s'intéresse à la recherche étymologique et combien de tâtonnements, d'erreurs même, seraient alors évités ! Mais quelle équipe aura assez d'abnégation et de courage pour se lancer jamais dans cette énorme entreprise ?}Puisse ce grand et beau dictionnaire un jour voir le jour !}
\cc{707–25}{}{Ce passage, sans correspondre parfaitement du point de vue métrique à celui des vers 667-686, présente néanmoins de troublantes ressemblances avec ce dernier. Se reporter à Lætitia\nc{\italique{The songs of Aristophanes}, pp. 421 sqq. Je reprends pour une bonne part sa colométrie.} pour les parallélismes qu'on peut faire avec le texte tel qu'il nous a été conservé.}
\cc{710–2}{ἀλλ᾽ οὖν ἥκεις τ᾿ ὅθεν ἥκεις.}{Je reproduis ci-dessous la \nieme{21} note d'Alphonse Willems. La correction qu'il y propose fut sauvée du néant de l'oubli et adoptée par la plupart des éditeurs modernes grâce à la diligente entremise de Coulon :

«~Ce passage a été refait de cinq ou six façons, plus bizarres les unes que les autres. 
La correction me paraît élémentaire. Les manuscrits concordent, sauf sur le premier ἥκεις (ἥκεις γ' dans Suidas, ἥκεις τ' dans le \italique{Ravennas}, ἥκισθ' dans le manuscrit de Munich [qui n'est qu'une copie du manuscrit de Ravenne, comportant quelques émendations parfois bien senties]). Lisez ἥκων γ', le sens devient parfait :
\citComm{ἀλλ᾽ οὖν ἥκων γ᾽ ὅθεν ἥκεις}
Il est à peine besoin d'ajouter que ἥκων est pris dans le sens de ἀνήκων, \italique{redux}, ce qui est des plus fréquents\nc{\italique{Paix}, 265, 275; \italique{Lysistr}., 595, 731, 924 ; \italique{Gren}., 503 ; \italique{Ploutos}, 521, 1179, etc.}. Et quant à ἀλλ᾽ οὖν..γε, chacun sait qu'il signifie \italique{toujours}, \italique{tout au moins}, \italique{en tous cas}\nc{\italique{Nuées}, 985, 1002 ; \italique{Guêpes}, 1129, 1190 ; Esch., \italique{Prom}., 1058 ; Soph., \italique{El}., 233 ; \italique{Ant}., 84, etc. M'éloignant de Willems quant au sens à donner à ce vers, je ne retranscris pas ici la traduction qu'il en propose. Le lecteur la trouvera en annexe où toutes les notes d'Alphonse Willems sur les \italique{Thesmophories} sont intégralement reproduites.}.\\
C'est le mot d'Épicharme, fr. 126\nc{$[=$ Plutarque, \italique{Moralia}, Παραμυθητικὸς πρὸς Ἀπολλώνιον, 110a $=$ fr. 213(245) Epich. Kassel-Austin, \italique{Poet. Com. Gr.}, I, \italique{Comoedia Dorica}.$]$} : κἀπῆνθεν ὅθεν ἦνθεν πάλιν ; ou de Ménandre (t. IV, p. 211 [$=$ fr. 373 K.-A.]) : ἀπῆλθεν ὅθεν ἦλθεν ταχύ. Cf. aussi Soph., \italique{Œd. Col.}, 1227 : βῆναι κεῖσ' ὁπόθεν περ ἥκει ; Eur., \italique{Her. F.}, 1247 : θανών, ὅθενπερ ἦλθον, εἶμι γῆς ὕπο, etc.~»

À la lumière de la démonstration de Willems, je me permets d'ajouter la glose suivante du \latin{Lexicon} de Photius (η 88) :\citComm{ἥκειν· ἀντὶ τοῦ ἥκων.}Supposer que la glose de Photius se rapporte à notre vers offre l'avantage de rester plus proche du texte des manuscrits : il y aurait eu simple permutation d'un ν par un ς.
L'on peut donc lire après Willems :\citComm{ἀλλ᾽ οὖν ἥκειν γ᾽ ὅθεν ἥκεις.\\
\francais{Au moins, de revenir d'où tu es parti.}}}
\cc{711–2}{φαύλως τ᾽ ἀποδρὰς οὐ λέξεις / οἷον δράσας διέδυς ἔργον.}{Où se trouve le complément réclamé par λέγειν ? Je ne suis pas sûr qu'il soit de bonne grécité de construire la phrase ainsi : λέξεις οἷον δράσας, «~tu raconteras tout ce que tu as fait~». (Sauf erreur de ma part, le grec aurait recouru à un verbe conjugué et non au participe de δράω.)
Une légère retouche du texte permettrait une bien meilleure lecture, tant pour la syntaxe grecque qu'au niveau du sens.
Lisez : 
\citComm{φαύλως γ᾽ ἀποδρὰς οὐκ ἔξει.\\
\francais{Ce n'est pas avec une fuite facile que tu t'en iras !}}
(Confusion entre un Κ et un Λ.)
Le vers 712, οἷον δράσας διέδυς ἔργον, est en fait apposé au vers 713 et en explique la sentence : λήψει δὲ κακόν, « pour tout ce que tu as fait, tu recevras châtiment ».
Pour justifier cette correction, je m'appuierai sur un débris de scholie qui se trouve à l'intérieur même du texte du \italique{Ravennas}, et que j'interprète ainsi :
\citComm{ἥκεις, φησίν, εἰς τὰς χεῖρας ἡμῶν, ὅθεν οὐ φεύξῃ.\\
Tu [re]viens, dit-elle, dans nos bras, et d'où tu ne fuiras pas.}
Comprise ainsi, cette scholie est clairement une paraphrase du vers 710 pour la première partie (ἥκεις εἰς τὰς χεῖρας ἡμῶν $=$ ἥκειν γ᾽ ὅθεν ἥκεις). C'est la seconde partie de la scholie qui nous intéresse : ὅθεν οὐ φεύξῃ, «~d'où tu ne fuiras pas~». Peut-on espérer meilleure glose que `οὐ φεύξῃ' pour le `οὐκ ἔξει' que je rétablis ?}
